% !TeX root = ../../Book3.tex
% 纲要标题：## 附录 A　代数数论初步：理想、整数环与分歧
% 纲要位置：工作记录/计划纲要.md:2816

\chapter{代数数论初步：理想、整数环与分歧}
\label{b3:app:A}

\begin{chapterabstract}
数域的整数环把整性变成算术上的整数概念。本附录证明它是有限秩整数格，进而由离散赋值环得到理想的唯一分解；再用二次数域计算整基、判别式、素理想分解与分歧。局部化和完备化说明如何单独考察一个素数处的行为，理想类群衡量理想能否由一个元素生成，Dedekind ζ 函数则用理想范数组织起所有非零理想。
\end{chapterabstract}

\begin{learninggoals}
\begin{itemize}
\item 用迹配对证明整数环的有限生成性，区分整数环与给定整生成元所生成的子环。
\item 从 DVR 的局部次数读取理想分解，解释 Dedekind--Kummer 条件 \(p\nmid[\mathcal O_K:\mathbb Z[\theta]]\) 为何必要。
\item 计算二次数域的整基、判别式、分歧指数和剩余次数。
\item 区分单个素理想处的完备化与 \(K\otimes_{\mathbb Q}\mathbb Q_p\) 的全部因子。
\item 检验一个理想非主生成，说明它在理想类群中给出的非零类，并核对 Dedekind ζ 函数的素理想乘积在实数 \(s>1\) 时的收敛条件。
\end{itemize}
\end{learninggoals}

在二次数域 \(\mathbb Q(\sqrt5)\) 中，\(\alpha=(1+\sqrt5)/2\) 不属于 \(\mathbb Z[\sqrt5]\)，却满足 \(\alpha^2-\alpha-1=0\)。若把“整数”仅理解为基 \(1,\sqrt5\) 下系数都是普通整数，就会漏掉这个元素。首一整系数方程给出更稳固的定义，并由此确定数域真正的整数环。

\ProseParagraphBreak

整数环作为 \(\mathbb Z\) 模形成有限秩格；选一组整基以后，迹配对的判别式可以检验局部素数处的行为。然而元素未必唯一分解为不可约元，理想却能在各素理想处记录赋值并唯一分解。把整基、判别式与理想分解放在同一个二次域里实际计算，才能看清这些数据各回答什么问题。

\ProseParagraphBreak

预备知识是\chapref{b3:ch:05}的整扩张、\chapref{b3:ch:10}的 DVR、\chapref{b3:ch:13}的完备化、\chapref{b3:ch:18}的正则局部环，以及第一卷第18章的迹、范数和判别式。整数格的证明还使用第二卷第4章中 PID 上有限生成模的结构定理。局部算例使用第23章的单根提升；最后一个级数只要求实数项级数的基本收敛判别。

\ProseParagraphBreak
第10章的 Dedekind 整环理论在这里成为素数分解的工具：各局部环中的统一化元幂次给出分歧指数，剩余域的大小给出剩余次数。整基与判别式将这些局部数据同整数格联系起来。Dedekind 整环的等价刻画与理想唯一分解还可参阅松村英之\cite[Theorem 11.6, pp. 82--84]{Matsumura1989CommutativeRingTheory}。

% !TeX root = ../../Book3.tex
% 纲要标题：### A.1 整闭包的算术意义
% 纲要位置：工作记录/计划纲要.md:2818

\section{整闭包的算术意义}
\label{b3:sec:A01}

% 纲要内容（仅作写作计划，尚非教材正文）：
%
% * 清分母恢复整个数域，具体首一方程的计算在正文完成
% * 迹的最小多项式公式、非退化迹矩阵与有限候选余类；Noether 性保证有限生成，PID 模结构给出自由性
% * 用 Lying-over 得到非零素理想，结合非零素理想均极大证明维数恰为一
% <!-- 短节：不设小节 -->
%
% ---
%


有理整数的整除问题进入数域以后，首先需要确定哪些元素应当算作这个域中的整数。只取某个生成元的整数多项式往往不够。例如在 \(\mathbb Q(\sqrt5)\) 中，\(\sqrt5\) 是整的，\((1+\sqrt5)/2\) 也满足首一整系数方程；环 \(\mathbb Z[\sqrt5]\) 却没有包含后一个元素。整闭包使这个选择不再依赖生成元。

\ProseParagraphBreak
沿用\cref{b3:def:number-field-integers}，数域是有限扩张 \(K/\mathbb Q\)，它的整数环为
\[
\mathcal O_K=\{a\in K:a\;\text{在}\;\mathbb Z\;\text{上整}\}.
\]
整数环的 Dedekind 性已在\cref{b3:thm:number-integers-dedekind}证明，使第10章的理想理论可以用于数域。本节给出另一个迹矩阵证明：它的算术用途在于把整数环夹在两个显式整数格之间，给整元素的坐标分母一个统一界，从而为寻找整基和计算判别式提供有限候选。

\ProseParagraphBreak
由\cref{b3:prop:integral-closure-ring}，这个集合确实是子环。这里的“整数”不要求实数意义下的大小或正负；它表达的是首一多项式关系。整数环也足够大：把非零普通整数用作分母，就能恢复整个数域。

\begin{lemma}\label{b3:lem:number-field-clear-denominators}
设 \(K/\mathbb Q\) 为有限扩张，\(\mathcal O_K\) 为 \(K\) 的代数整数环，\(a\in K\)。存在非零整数 \(d\)，使 \(da\in\mathcal O_K\)。特别地，\(\mathbb Q\mathcal O_K=\operatorname{Frac}(\mathcal O_K)=K\)。若 \(a^3+a/2+1/3=0\)，则 \(b=6a\) 满足首一整系数方程 \(b^3+18b+72=0\)。
\end{lemma}
\begin{proof}
写 \(a^r+c_1a^{r-1}+\cdots+c_r=0\)，其中 \(c_i\in\mathbb Q\)。取正整数 \(d\) 使每个 \(dc_i\) 都为整数，则
\[
(da)^r+(dc_1)(da)^{r-1}+(d^2c_2)(da)^{r-2}+\cdots+d^rc_r=0
\]
是首一整系数关系。最后 \(a=(da)/d\)，即得有理张成与分式域的结论。对陈述中的三次方程，取 \(d=6\)，代入 \(b=6a\) 并乘以 \(6^3\)，就得到 \(b^3+18b+72=0\)。这只证明 \(6a\) 整；整元素除以普通整数未必仍整，例如 \(1/2\) 不是代数整数。
\end{proof}

下面先证明整元素的迹为整数，再用非退化迹矩阵控制坐标分母。有限生成性与自由性由格包含得到；最后将这些模论结论转成 Noether 性、整闭性和维数一。

\begin{theorem}\label{b3:thm:number-field-integral-lattice}
设 \(K/\mathbb Q\) 为次数 \(n\) 的扩张，\(\mathcal O_K\) 为 \(K\) 的代数整数环，\(\beta_1,\ldots,\beta_n\) 为由整元素组成的 \(\mathbb Q\)-基，\(L=\sum_i\mathbb Z\beta_i\)。令
\(D=\det(\operatorname{Tr}_{K/\mathbb Q}(\beta_i\beta_j))\)。
则 \(D\in\mathbb Z\setminus\{0\}\)，且
\[
L\subseteq\mathcal O_K\subseteq D^{-1}L.
\]
特别地，\(\mathcal O_K\) 是秩 \(n\) 的自由 \(\mathbb Z\)-模，也是一维 Noether 整闭整环。
\end{theorem}
\begin{proof}
这样的整元素基可由任意有理基分别清分母得到。若 \(u\in\mathcal O_K\)，它的每个 \(\mathbb Q\)-共轭也满足同一个首一整系数方程，因而仍整。最小多项式的系数是这些共轭根的初等对称多项式，所以是代数整数；另一方面，它们属于 \(\mathbb Q\)。有理数中的代数整数恰为普通整数，故最小多项式属于 \(\mathbb Z[T]\)。写为 \(m_u(T)=T^r+c_1T^{r-1}+\cdots+c_r\)，则由可分扩张的迹及其传递公式，
\[
\operatorname{Tr}_{\mathbb Q(u)/\mathbb Q}(u)=-c_1,\qquad
\operatorname{Tr}_{K/\mathbb Q}(u)=[K:\mathbb Q(u)](-c_1)\in\mathbb Z.
\]

特征零扩张可分，其迹配对非退化，见第一卷定理18.3.3。故整数矩阵
\[
T=\bigl(\operatorname{Tr}_{K/\mathbb Q}(\beta_i\beta_j)\bigr)_{i,j}
\]
的行列式 \(D\) 非零。对任意 \(u=\sum_jc_j\beta_j\in\mathcal O_K\)，各 \(\operatorname{Tr}(\beta_i u)\) 均为整数，因而 \(T(c_j)_j\in\mathbb Z^n\)。乘以伴随矩阵得 \(Dc_j\in\mathbb Z\)，所以
\[
L\subseteq\mathcal O_K\subseteq D^{-1}L.
\]
右端 \(D^{-1}L\cong\mathbb Z^n\) 是有限生成 \(\mathbb Z\)-模。由于 \(\mathbb Z\) 是 Noether 环，其子模 \(\mathcal O_K\) 也有限生成。又因 \(\mathcal O_K\subseteq K\)，它作为 \(\mathbb Z\)-模无挠；\(\mathbb Z\) 是主理想整环，第二卷定理4.3.2的模结构定理因此给出 \(\mathcal O_K\) 自由。两侧的 \(L\) 与 \(D^{-1}L\) 都张成 \(n\) 维 \(\mathbb Q\)-空间 \(K\)，所以 \(\mathcal O_K\) 的秩恰为 \(n\)。

\(\mathcal O_K\) 的每个理想都是这个有限生成 \(\mathbb Z\)-模的子模，故有有限组 \(\mathbb Z\)-模生成元；这组元素同时生成该理想。因此 \(\mathcal O_K\) 是 Noether 环。若 \(x\in K\) 在 \(\mathcal O_K\) 上整，整性的传递性使它在 \(\mathbb Z\) 上整，故 \(x\in\mathcal O_K\)。

最后考察维数。设 \(\mathfrak p\) 为非零素理想，取 \(0\ne a\in\mathfrak p\)。其最小多项式的常数项 \(c\) 非零，并且多项式关系给出 \(c\in a\mathcal O_K\subseteq\mathfrak p\)。因此 \(\mathfrak p\cap\mathbb Z\) 是非零素理想，必为 \((p)\)，其中 \(p\) 是素数。

环 \(\mathcal O_K/\mathfrak p\) 是有限维 \(\mathbb F_p\)-空间，又是整环，故为有限域。因此 \(\mathfrak p\) 极大，所有非零素理想都不能再严格包含于另一个真素理想。由于 \(\mathcal O_K\) 是整环，每条严格素理想链的长度至多为一。

反过来，任取素整数 \(p\)，对整扩张 \(\mathbb Z\subseteq\mathcal O_K\) 应用\cref{b3:thm:lying-over}，得到素理想 \(\mathfrak q\) 满足 \(\mathfrak q\cap\mathbb Z=(p)\)。因 \(p\in\mathfrak q\)，\(\mathfrak q\ne(0)\)，于是 \((0)\subsetneq\mathfrak q\) 是长度一的素理想链。因此维数恰为一。
\end{proof}

这个证明给出的不仅是有限生成性。固定基 \(\beta_i\) 后，每个整元素的坐标 \(c_i\) 都满足 \(Dc_i\in\mathbb Z\)，其分母统一受到 \(D\) 控制。有限群 \(D^{-1}L/L\) 因而包含全部候选余类；同一余类中的元素相差整元素，所以只须对每个余类选一个代表检验整性。这提供了寻找整基的有限性依据，却没有断言任意整元素生成的子环已经等于整数环。

\ProseParagraphBreak
\cref{b3:def:dedekind-domain}已经定义 Dedekind 整环，并允许域作为零维情形；非域情形恰为一维 Noether 整闭整环。整数环属于后一种情形。由\cref{b3:thm:dvr-characterizations}，在它的任意非零素理想处局部化，得到离散赋值环。整数环未必有元素的唯一分解，然而每个这样的局部环都可选取统一化元，以其幂次衡量整除次数。下一节将把这些局部次数用于理想的唯一分解。

% !TeX root = ../../Book3.tex
% 纲要标题：### A.2 数域整数环、判别式与理想分解
% 纲要位置：工作记录/计划纲要.md:2785

\section{数域整数环、判别式与理想分解}
\label{b3:sec:A02}

% 纲要内容（仅作写作计划，尚非教材正文）：
% * 复用第10章的理想唯一分解，说明局部指数及有限商环
% * 整基、域判别式、子环指数平方公式与二次数域整数环
% * 理想范数的乘法性、剩余次数与分歧指数；Dedekind--Kummer 的指数互素条件及局部长度证明
% * 惰性、分歧与完全分裂的区别；判别式检测分歧
% #### A.2.1 理想唯一分解
% #### A.2.2 整基与判别式
% #### A.2.3 素数分解与 Dedekind--Kummer 定理
% #### A.2.4 判别式与分歧

在 \(\mathbb Q(\sqrt5)\) 的整数环中，有理素数 \(2\) 仍生成一个素理想，\(5\) 生成的理想却是一个素理想的平方，而 \(11\) 生成的理想分成两个不同的素理想。要区别这些情形，单看整数的因子分解已不够：必须计算素理想的幂指数以及相应剩余域的大小。整数环的一组整基让这些局部数据能够与整个数域的次数相比较，判别式则检测哪些素数出现重复的素理想因子。

\subsection{理想唯一分解}
\label{b3:subsec:A02:01}

\begin{theorem}\label{b3:thm:integer-ideal-factorization}
设 \(K\) 为数域，\(\mathcal O_K\) 为其代数整数环，\(R=\mathcal O_K\)。每个非零真理想 \(I\subset R\) 唯一写成
\[
I=\mathfrak p_1^{a_1}\cdots\mathfrak p_r^{a_r},
\qquad a_i\ge1,
\]
其中 \(\mathfrak p_i\) 为互不相同的非零素理想。
\end{theorem}
\begin{proof}
这是\cref{b3:thm:dedekind-factorization}对整数环的应用。为把指数与算术联系起来，先注意上一节的常数项论证说明 \(I\) 含有某个非零整数 \(c\)。因 \(R\) 是有限秩整数格，\(R/I\) 有限，故只有有限多个极大理想包含 \(I\)，记为 \(\mathfrak p_1,\ldots,\mathfrak p_r\)。指数由
\[
IR_{\mathfrak p_i}=(\mathfrak p_iR_{\mathfrak p_i})^{a_i}
\]
唯一确定。第10章已经证明，这些局部等式通过局部检测给出所列全局唯一分解。
\end{proof}

一般等价条件可查 Matsumura 的分解定理\cite[Theorem 11.6, pp. 82--84]{Matsumura1989CommutativeRingTheory}。在数域整数环中，商环还是有限集合，因此理想分解也能给出可计算的整数。

\subsection{整基与判别式}
\label{b3:subsec:A02:02}

\begin{definition}\label{b3:def:integer-discriminant-norm}
设 \(K/\mathbb Q\) 的次数为 \(n\)，\(\mathcal O_K\) 为 \(K\) 的代数整数环。若 \(\omega_1,\ldots,\omega_n\) 是 \(\mathcal O_K\) 的一组 \(\mathbb Z\)-基，则称之为一组\term[integralbasis]{整基}[integral basis]。由这组整基定义的整数
\[
d_K=\det\bigl(\operatorname{Tr}_{K/\mathbb Q}(\omega_i\omega_j)\bigr)
\]
称为 \(K\) 的判别式。
\end{definition}

两组整基之间的变换矩阵行列式为 \(\pm1\)，迹矩阵的行列式乘以该行列式的平方，所以 \(d_K\) 与整基无关。更一般地，若子环 \(S\subseteq\mathcal O_K\) 也是秩 \(n\) 的整数格，则
\[
\operatorname{disc}(S)=[\mathcal O_K:S]^2d_K.
\]
这是整数矩阵的指数公式与迹矩阵的换基公式的直接结合。判别式可能带有来自子环指数的额外平方因子，不能在未检查整数环时就把它们全都解释为分歧。

\begin{example}\label{b3:exm:quadratic-integer-ring}
设 \(d\ne0,1\) 为平方自由整数，\(K=\mathbb Q(\sqrt d)\)。则
\[
\mathcal O_K=
\begin{cases}
\mathbb Z[(1+\sqrt d)/2],&d\equiv1\pmod4,\\
\mathbb Z[\sqrt d],&d\equiv2,3\pmod4,
\end{cases}
\qquad
d_K=
\begin{cases}
d,&d\equiv1\pmod4,\\
4d,&d\equiv2,3\pmod4.
\end{cases}
\]
为证明整数环公式，写一个整元素为 \(r+s\sqrt d\)，其中 \(r,s\in\mathbb Q\)。迹和范数给出
\(a=2r\in\mathbb Z\) 以及 \(r^2-ds^2\in\mathbb Z\)。令 \(b=2s\in\mathbb Q\)，则 \(db^2\in\mathbb Z\)。若 \(b=u/v\) 为既约分数，便有 \(v^2\mid d\)，平方自由性迫使 \(v=1\)。于是
\[
a,b\in\mathbb Z,\qquad a^2-db^2\equiv0\pmod4.
\]
若 \(d\equiv2,3\pmod4\)，\(b\) 不可能为奇数，进而 \(a,b\) 都为偶数；若 \(d\equiv1\pmod4\)，条件等价于 \(a,b\) 同奇偶。反过来，这些条件使迹、范数均为整数，所以相应元素满足首一二次整系数方程。最后在所列整基中直接计算迹矩阵，即得判别式。
\end{example}

\subsection{素数分解与 Dedekind--Kummer 定理}
\label{b3:subsec:A02:03}

\begin{definition}\label{b3:def:integer-ideal-norm}
设 \(K\) 为数域，\(\mathcal O_K\) 为其代数整数环，\(I\subseteq\mathcal O_K\) 为非零理想。有限商环 \(\mathcal O_K/I\) 的基数
\[
N(I)=|\mathcal O_K/I|
\]
称为 \(I\) 的\term[idealnorm]{理想范数}[ideal norm]。其中 \(N(\mathcal O_K)=1\)。
\end{definition}

\begin{proposition}\label{b3:prop:prime-decomposition-degree}
设 \(K/\mathbb Q\) 的次数为 \(n\)，\(\mathcal O_K\) 为 \(K\) 的代数整数环，\(p\) 为有理素数，并写
\[
p\mathcal O_K=\prod_{i=1}^r\mathfrak p_i^{e_i},
\]
其中 \(\mathfrak p_i\) 为互不相同的非零素理想，\(e_i\) 为正整数。则存在正整数 \(f_i\)，使 \(\mathcal O_K/\mathfrak p_i\cong\mathbb F_{p^{f_i}}\)，且对每个非负整数 \(a\)，有
\[
\sum_{i=1}^r e_if_i=n,\qquad
N(\mathfrak p_i^a)=p^{af_i}.
\]
称 \(e_i\) 为 \(\mathfrak p_i\) 在 \(p\) 上方的分歧指数，称 \(f_i\) 为其剩余次数。
\end{proposition}
\begin{proof}
记 \(R=\mathcal O_K\)。每个剩余域 \(R/\mathfrak p_i\) 都是有限商环 \(R/pR\) 的商，因而是某个有限 \(\mathbb F_p\)-扩张。固定 \(i\)，对每个非负整数 \(j\)，令
\[
V_j=\mathfrak p_i^j/\mathfrak p_i^{j+1}.
\]
由于 \(\mathfrak p_iV_j=0\)，\(V_j\) 本来就是 \(R/\mathfrak p_i\)-向量空间。任意 \(s\notin\mathfrak p_i\) 在剩余域中可逆，因此局部化映射 \(V_j\rightarrow(V_j)_{\mathfrak p_i}\) 是同构。若 \(\pi_i\) 为 DVR \(R_{\mathfrak p_i}\) 的统一化元，则
\[
(V_j)_{\mathfrak p_i}
\cong\pi_i^jR_{\mathfrak p_i}/\pi_i^{j+1}R_{\mathfrak p_i}
\cong R_{\mathfrak p_i}/\mathfrak p_iR_{\mathfrak p_i}
\cong R/\mathfrak p_i.
\]
所以 \(V_j\) 的剩余域维数为一，基数为 \(p^{f_i}\)。对 \(R/\mathfrak p_i^a\) 使用由理想幂给出的有限滤过，逐层计算基数，得到 \(N(\mathfrak p_i^a)=p^{af_i}\)；\(a=0\) 时公式也成立。不同极大理想的幂两两互素，中国剩余定理及 \(R/pR\) 的 \(\mathbb F_p\)-维数 \(n\) 给出
\[
p^n=|R/pR|=\prod_i|R/\mathfrak p_i^{e_i}|=p^{\sum_i e_if_i}.
\]
\end{proof}

同样由中国剩余定理，若非零理想 \(I=\prod_i\mathfrak p_i^{a_i}\)，则 \(N(I)=\prod_iN(\mathfrak p_i)^{a_i}\)。理想相乘使各素理想指数相加，因此 \(N(IJ)=N(I)N(J)\)。

\begin{definition}\label{b3:def:numberfield-prime-splitting}
设 \(K/\mathbb Q\) 的次数为 \(n\)，\(\mathcal O_K\) 为其代数整数环，\(p\) 为有理素数，并写
\[
p\mathcal O_K=\prod_{i=1}^r\mathfrak p_i^{e_i},\qquad
f_i=[\mathcal O_K/\mathfrak p_i:\mathbb F_p],
\]
其中各 \(\mathfrak p_i\) 互不相同。若某个 \(e_i>1\)，则称 \(p\) 在 \(K\) 中\term[fenqi]{分歧}[ramified]；若所有 \(e_i=1\)，则称其未分歧。若 \(r=1\) 且 \(e_1=1\)（此时 \(f_1=n\)），则称 \(p\) 在 \(K\) 中\term[duoxing]{惰性}[inert]。若 \(r=n\) 且所有 \(e_i=f_i=1\)，则称 \(p\) 在 \(K\) 中完全分裂。
\end{definition}

要计算上述指数和剩余次数，可以将整数环的关系式模 \(p\) 约化。若一个整元素 \(\theta\) 生成的子环 \(\mathbb Z[\theta]\) 在 \(\mathcal O_K\) 中的指数不被 \(p\) 整除，这次约化就不会丢失整数环模 \(p\) 的信息。下面的 Dedekind--Kummer 定理把多项式的因子分解转化为素理想分解。\footnote{库默（Ernst Eduard Kummer，1810-1893），德国数学家，创立理想数理论，并利用它研究 Fermat 方程。}

\begin{proposition}[Dedekind--Kummer]\label{b3:prop:monogenic-prime-factorization}
设 \(K=\mathbb Q(\theta)\)，\(\theta\) 为整元素，\(\mathcal O_K\) 为 \(K\) 的代数整数环，\(f\in\mathbb Z[T]\) 为 \(\theta\) 的首一最小多项式。设素数 \(p\) 不整除 \([\mathcal O_K:\mathbb Z[\theta]]\)，并在 \(\mathbb F_p[T]\) 中分解
\[
\overline f=\prod_{i=1}^r g_i^{e_i},
\]
其中 \(g_i\) 为互不相同的首一不可约多项式。取任意整系数提升 \(\widetilde g_i\)，则
\[
p\mathcal O_K=\prod_{i=1}^r
\bigl(p,\widetilde g_i(\theta)\bigr)^{e_i},
\]
各括号内的理想为素理想，其剩余次数为 \(\deg g_i\)。
\end{proposition}
\begin{proof}
令 \(S=\mathbb Z[\theta]\)、\(R=\mathcal O_K\)。有限群 \(R/S\) 的阶与 \(p\) 互素，乘以 \(p\) 在此群上是同构。因此 \(S\cap pR=pS\) 且 \(S+pR=R\)，给出
\[
R/pR\simeq S/pS\simeq\mathbb F_p[T]/(\overline f).
\]
右端的极大理想对应各 \(g_i\)，剩余域次数为 \(\deg g_i\)，故 \(R\) 中位于 \(p\) 上方的素理想恰为 \(\mathfrak p_i=(p,\widetilde g_i(\theta))\)。在相应极大理想处局部化时，其余 \(g_j\) 都变为单位，因此有局部环同构
\[
(R/pR)_{\mathfrak p_i}
\cong R_{\mathfrak p_i}/pR_{\mathfrak p_i}
\cong\mathbb F_p[T]_{(g_i)}/(g_i^{e_i}).
\]
由理想唯一分解，存在正整数 \(a_i\)，使
\[
pR_{\mathfrak p_i}
=(\mathfrak p_iR_{\mathfrak p_i})^{a_i}.
\]
DVR 商 \(R_{\mathfrak p_i}/pR_{\mathfrak p_i}\) 按其极大理想的幂滤过，有 \(a_i\) 个非零层，各层均为其剩余域上的一维空间，所以合成长度为 \(a_i\)。右端按 \(g_i\) 的幂滤过，有 \(e_i\) 个非零层，各层均同构于 \(\mathbb F_p[T]/(g_i)\)，所以合成长度为 \(e_i\)。局部环同构保持合成长度，故 \(a_i=e_i\)，得到所述分解。
\end{proof}

\begin{example}\label{b3:exm:quadratic-five-splitting}
在 \(K=\mathbb Q(\sqrt5)\) 中取 \(\omega=(1+\sqrt5)/2\)，则
\(\mathcal O_K=\mathbb Z[\omega]\)，最小多项式为 \(T^2-T-1\)。模 \(2\) 不可约，故 \(2\mathcal O_K\) 本身为素理想，剩余域为 \(\mathbb F_4\)，其分歧指数为 \(1\)、剩余次数为 \(2\)，所以 \(2\) 惰性。模 \(5\) 等于 \((T-3)^2\)，故
\[
5\mathcal O_K=(5,\omega-3)^2,
\]
这里分歧指数为 \(2\)、剩余次数为 \(1\)。模 \(11\) 的两个根为 \(4,8\)，所以
\[
11\mathcal O_K=(11,\omega-4)(11,\omega-8),
\]
两个素理想的分歧指数和剩余次数都为 \(1\)，所以 \(11\) 完全分裂。三个素数分别展示惰性、分歧和完全分裂。若改用 \(\mathbb Z[\sqrt5]\)，它的指数为 \(2\)、判别式为 \(20\)；模 \(2\) 的重因子不能用于以上判据，因为指数条件恰好在这里失败。
\end{example}

\subsection{判别式与分歧}
\label{b3:subsec:A02:04}

\begin{proposition}\label{b3:prop:discriminant-ramification}
设 \(K\) 为数域，\(d_K\) 为其整数环 \(\mathcal O_K\) 的判别式，\(p\) 为有理素数。则 \(p\) 在 \(K\) 中分歧，当且仅当 \(p\mid d_K\)。
\end{proposition}
\begin{proof}
整基使 \(R=\mathcal O_K\) 上的乘法矩阵约化为 \(R/pR\) 上的乘法矩阵。因此 \(d_K\) 模 \(p\) 是有限维 \(\mathbb F_p\)-代数 \(C=R/pR\) 的迹配对的行列式。
若 \(C\) 有非零幂零元 \(a\)，则每个 \(ab\) 都幂零，乘法算子也幂零，故 \(\operatorname{Tr}(ab)=0\)；迹配对退化。若 \(C\) 约化，有限维交换约化代数为有限域扩张之积，而有限域上的有限扩张可分，各迹配对均非退化。因此配对非退化恰好等价于 \(C\) 约化。由理想分解及 DVR 中的商环计算，\(C\) 约化恰好等价于每个 \(e_i=1\)。
\end{proof}

% !TeX root = ../../Book3.tex
% 纲要标题：### A.3 素数处的局部化与完备化
% 纲要位置：工作记录/计划纲要.md:2839

\section{素数处的局部化与完备化}
\label{b3:sec:A03}

% 纲要内容（仅作写作计划，尚非教材正文）：
%
% * 统一化元值为一的归一化赋值；局部化、逆极限完备化与赋值度量完备域的识别
% * 张量表示、中国剩余定理和共尾理想幂给出全部局部因子；模 p 商的滤过证明各因子的秩
% * 分解指数、剩余次数、理想范数、完备域次数及判别式的全局与局部联系
% * 二次数域 Q(√5) 的模素数因子、e/f、范数、判别式和完备因子对照；模121数位与具体嵌入
% <!-- 短节：不设小节 -->
%
% ---
%


固定数域整数环 \(R=\mathcal O_K\) 中的非零素理想 \(\mathfrak p\)。局部化 \(R_{\mathfrak p}\) 允许把所有不属于 \(\mathfrak p\) 的元素放到分母中。因此它保留的是在这个素理想处的整除次数；一个元素在别的素理想处是否有分母，不再妨碍它属于 \(R_{\mathfrak p}\)。

\begin{definition}\label{b3:def:number-field-prime-completion}
设 \(K\) 为数域，\(\mathcal O_K\) 为其代数整数环，\(\mathfrak p\subset\mathcal O_K\) 为非零素理想。记 \(v_{\mathfrak p}:K^\times\to\mathbb Z\) 为离散赋值环 \((\mathcal O_K)_{\mathfrak p}\) 的归一化离散赋值，即对任一统一化元 \(\pi_{\mathfrak p}\) 有 \(v_{\mathfrak p}(\pi_{\mathfrak p})=1\)。记
\[
\widehat{\mathcal O}_{K,\mathfrak p}
=\varprojlim_r(\mathcal O_K)_{\mathfrak p}/
(\mathfrak p(\mathcal O_K)_{\mathfrak p})^r,
\qquad
K_{\mathfrak p}=\operatorname{Frac}(\widehat{\mathcal O}_{K,\mathfrak p}).
\]
分别称它们为相应的完备整数环与 \(K\) 在 \(\mathfrak p\) 处的完备化。
\end{definition}
\symidx{\(\widehat{\mathcal O}_{K,\mathfrak p}\)、\(K_{\mathfrak p}\)：代数整数环的局部完备化及其分式域}

由\cref{b3:thm:dvr-characterizations,b3:thm:regular-local-completion}，\(\widehat{\mathcal O}_{K,\mathfrak p}\) 仍为 DVR。完备化保留极大理想的生成元与剩余域，所以原来的统一化元仍为统一化元。原环嵌入完备化由\cref{b3:thm:krull-intersection}保证。赋值也因而延拓到 \(K_{\mathfrak p}\)。

这个代数定义也给出通常的赋值度量完备化。令 \(V=R_{\mathfrak p}\)，取统一化元 \(\pi\)，并约定 \(v_{\mathfrak p}(0)=\infty\)。度量 \(d(x,y)=2^{-v_{\mathfrak p}(x-y)}\) 的 Cauchy 条件，就是对每个 \(r\)，充分靠后的差都属于 \(\pi^rV\)。对 \(V\) 中的序列，这等价于每个模 \(\pi^r\) 余类最终稳定；相容的稳定余类恰为 \(\widehat V\) 的元素。反过来，逐阶选取这些余类的代表便得到 Cauchy 序列。因此 \(\widehat V\) 是 \(V\) 在此度量下的完备化。

在 \(K\) 中，Cauchy 序列的赋值最终有统一下界：固定一个足够靠后的项，再用差的赋值非负即可。乘以某个固定的 \(\pi^N\) 后，整个尾部便落在 \(V\) 中，故极限落在 \(\pi^{-N}\widehat V\)。另一方面，完备 DVR 的每个分式都属于某个 \(\pi^{-N}\widehat V\)，并能由 \(K\) 中的序列逼近。由此 \(\operatorname{Frac}(\widehat V)\) 正是 \(K\) 的赋值度量完备化；这个识别在 \(K\) 上为恒等，即保留 \(K\) 的原有嵌入。

\ProseParagraphBreak
局部化和完备化不是同一操作。例如 \(\mathbb Z_{(p)}\) 只由分母不被 \(p\) 整除的有理数组成，是可数集；\(\mathbb Z_p\) 则包含任意相容的 \(p\) 进数位串，是不可数集。二者有相同的模 \(p^r\) 商，却并不相等。这正是完备化补入相容极限的含义。

整体的 \(p\)-进完备化同时保留 \(p\) 上方的全部素理想。下面将有限自由模的张量积表示、中国剩余分解和共尾理想幂依次结合，得到各局部完备因子的乘积。

\begin{theorem}\label{b3:thm:number-field-completion-product}
设 \(K\) 为数域，\(R=\mathcal O_K\) 为其代数整数环，\(p\) 为有理素数，且
\(pR=\prod_{i=1}^r\mathfrak p_i^{e_i}\)。则自然映射给出
\[
R\otimes_{\mathbb Z}\mathbb Z_p
\simeq\prod_{i=1}^r\widehat{\mathcal O}_{K,\mathfrak p_i},
\qquad
K\otimes_{\mathbb Q}\mathbb Q_p
\simeq\prod_{i=1}^rK_{\mathfrak p_i}.
\]
若 \(f_i=[R/\mathfrak p_i:\mathbb F_p]\)，则
\([K_{\mathfrak p_i}:\mathbb Q_p]=e_if_i\)。
\end{theorem}
\begin{proof}
\(R\) 是有限自由 \(\mathbb Z\)-模，有限生成模的完备化与张量积比较给出
\(R\otimes\mathbb Z_p\simeq\varprojlim_a R/p^aR\)。中国剩余定理给出
\[
R/p^aR\simeq\prod_i R/\mathfrak p_i^{ae_i}.
\]
环 \(R/\mathfrak p_i^b\) 只有一个极大理想；不在 \(\mathfrak p_i\) 中的元素在这个商中已经可逆。因此它与
\(R_{\mathfrak p_i}/\mathfrak p_i^bR_{\mathfrak p_i}\) 相同。对任意 \(b\ge1\)，取 \(a\ge\lceil b/e_i\rceil\) 就有 \(ae_i\ge b\)。因此这些指数在所有正整数中共尾，两套理想幂定义同一拓扑和完备化。这是\secref{b3:sec:1301}中共尾滤过给出同一相容余类模的同一理由；取有限乘积的逆极限即得第一式。

在每个因子中 \(p\) 等于单位乘以统一化元的 \(e_i\) 次幂，故把 \(p\) 变成可逆元就得到该 DVR 的分式域。对第一式把 \(p\) 变成可逆元，便得第二式。

每个完备整数环是有限自由 \(\mathbb Z_p\)-模：它是有限自由模 \(R\otimes\mathbb Z_p\) 的直和因子，因而有限且无挠，在主理想整环 \(\mathbb Z_p\) 上自由。因为 \(p\) 是单位乘以 \(\pi_i^{e_i}\)，完备化保持有限阶商，故
\[
\widehat{\mathcal O}_{K,\mathfrak p_i}/p\widehat{\mathcal O}_{K,\mathfrak p_i}
\cong R_{\mathfrak p_i}/pR_{\mathfrak p_i}
\cong R/\mathfrak p_i^{e_i}.
\]
最后的商按极大理想幂有 \(e_i\) 层，每层的 \(\mathbb F_p\)-维数为 \(f_i\)，故其维数为 \(e_if_i\)。这也是该自由 \(\mathbb Z_p\)-模的秩。对 \(\widehat{\mathcal O}_{K,\mathfrak p_i}\) 倒置 \(p\)，得到其分式域 \(K_{\mathfrak p_i}\)；原来的自由模基在沿 \(\mathbb Z_p\to\mathbb Q_p\) 扩张标量后成为 \(\mathbb Q_p\) 基，所以 \([K_{\mathfrak p_i}:\mathbb Q_p]=e_if_i\)。
\end{proof}

此定理解释了为什么在一个有理素数处研究数域时，可能同时出现几个局部域。张量积保留 \(p\) 上方的全部素理想；选定 \(\mathfrak p_i\)，才是选定一个因子。因此 \(K\otimes_{\mathbb Q}\mathbb Q_p\) 一般不是域，虽然每个 \(K_{\mathfrak p_i}\) 都是域。

\begin{table}[htbp]
\centering
\caption[数域整数环的全局与局部数据]{数域整数环的全局与局部数据。这里 \(R=\mathcal O_K\)，\(n=[K:\mathbb Q]\)，\(p\) 为有理素数，\(\mathfrak p\mid p\) 表示 \(\mathfrak p\) 位于 \(p\) 上方；\(e_{\mathfrak p}\)、\(f_{\mathfrak p}\) 分别为分歧指数、剩余次数，\(d_K\) 为数域判别式。}
\label{b3:tab:number-field-global-local-dictionary}
\begin{tabularx}{\linewidth}{@{}>{\raggedright\arraybackslash}p{0.22\linewidth}X@{}}
\toprule
算术数据 & 全局与局部的联系 \\
\midrule
局部整数环 & \(R_{\mathfrak p}\) 倒置 \(\mathfrak p\) 外的元素，是剩余域为 \(R/\mathfrak p\) 的 DVR，通常尚未完备。 \\
完备整数环 & \(\widehat{\mathcal O}_{K,\mathfrak p}\) 加入各阶相容余类的极限，仍为 DVR，剩余域不变。 \\
素理想分解 & \(pR=\prod_{\mathfrak p\mid p}\mathfrak p^{e_{\mathfrak p}}\)，且 \(v_{\mathfrak p}(p)=e_{\mathfrak p}\)。 \\
剩余次数 & \(f_{\mathfrak p}=[\kappa(\mathfrak p):\mathbb F_p]\)；完备化的剩余域仍为 \(\kappa(\mathfrak p)=R/\mathfrak p\)。 \\
素理想范数 & \(N(\mathfrak p)=p^{f_{\mathfrak p}}\)。 \\
完备因子的次数 & \([K_{\mathfrak p}:\mathbb Q_p]=e_{\mathfrak p}f_{\mathfrak p}\)。 \\
全部因子的总次数 & \(\sum_{\mathfrak p\mid p}e_{\mathfrak p}f_{\mathfrak p}=n\)。 \\
判别式与分歧 & \(p\mid d_K\Longleftrightarrow\) 存在 \(\mathfrak p\mid p\) 满足 \(e_{\mathfrak p}>1\)。 \\
\bottomrule
\end{tabularx}
\end{table}

\cref{b3:tab:number-field-global-local-dictionary}中的赋值均按统一化元的值为一归一化。理想分解给出每个局部因子中的整除次数，剩余次数记录剩余域的大小，两者的乘积才是完备因子的扩张次数。

\begin{example}\label{b3:exm:quadratic-five-completions}\label{b3:exm:eleven-adic-lifts}
取 \(K=\mathbb Q(\sqrt5)\)、\(\omega=(1+\sqrt5)/2\)，其最小多项式为 \(f(T)=T^2-T-1\)。模 \(11\) 的根 \(4,8\) 给出两个素理想，其 \(e=f=1\)，所以上一定理先给出 \(K\otimes\mathbb Q_{11}\cong\mathbb Q_{11}\times\mathbb Q_{11}\)。Hensel 提升再把这两个因子的投影写成具体嵌入。

因 \(f'(4)=7\)、\(f'(8)=15\) 模 \(11\) 均非零，两根在 \(\mathbb Z_{11}\) 中各自唯一提升。为计算下一位，写为 \(4+11c\)、\(8+11d\)。模 \(121\) 分别得到
\[
11+77c\equiv0,\qquad55+165d\equiv0\pmod{121}.
\]
除以 \(11\) 后在 \(\mathbb F_{11}\) 中解得 \(c=3\)、\(d=7\)，即两根分别为 \(37,85\pmod{121}\)。若其完整提升为 \(a_1,a_2\)，令 \(\omega\mapsto a_i\) 就给出两个 \(K\longrightarrow\mathbb Q_{11}\) 的嵌入；\(\sqrt5=2\omega-1\) 的像互为相反数。

模 \(2\) 不可约，所得因子是二次未分歧扩张；模 \(5\) 则只有一个素理想，所得因子为二次分歧扩张。在后者中
\[
2v_{\mathfrak p}(\sqrt5)=v_{\mathfrak p}(5)=2,
\]
所以 \(\sqrt5\) 的归一化赋值为 \(1\)，是一个统一化元。三个素数处的数据汇于\cref{b3:tab:quadratic-five-local-data}。
\end{example}

\begin{table}[htbp]
\centering
\caption[\(\mathbb Q(\sqrt5)\) 在三个有理素数处的局部数据]{\(\mathbb Q(\sqrt5)\) 在三个有理素数处的局部数据。取 \(f(T)=T^2-T-1\)，域判别式 \(d_K=5\)；\((e,f)\) 与范数按各素理想分别列出。}
\label{b3:tab:quadratic-five-local-data}
\begin{tabularx}{\linewidth}{@{}c>{\raggedright\arraybackslash}p{0.23\linewidth}cc>{\raggedright\arraybackslash}X@{}}
\toprule
\(p\) & \(\overline f\) & \((e,f)\) & \(N(\mathfrak p)\) & 分解类型与 \(K\otimes_{\mathbb Q}\mathbb Q_p\) \\
\midrule
\(2\) & \(T^2+T+1\)，不可约 & \((1,2)\) & \(4\) & 惰性，二次未分歧域 \(\mathbb Q_2(\sqrt5)\) \\
\(5\) & \((T-3)^2\) & \((2,1)\) & \(5\) & 分歧，二次分歧域 \(\mathbb Q_5(\sqrt5)\) \\
\(11\) & \((T-4)(T-8)\) & 各 \((1,1)\) & 各 \(11\) & 完全分裂，\(\mathbb Q_{11}\times\mathbb Q_{11}\) \\
\bottomrule
\end{tabularx}
\end{table}

这里的单根提升可直接使用\cref{b3:thm:hensel-simple-root}。它帮助计算已经存在的完备因子；各因子的数目与次数，则先由理想分解确定。这样就把第10章的赋值、第13章的完备化与第23章的 Hensel 提升用于同一个算例。

% !TeX root = ../../Book3.tex
% 纲要标题：### A.4 理想类群与 Dedekind ζ 函数
% 纲要位置：工作记录/计划纲要.md:2806

\section{理想类群与 Dedekind \texorpdfstring{\(\zeta\)}{\AlgebraPdfZeta} 函数}
\label{b3:sec:A04}

% 纲要内容（仅作写作计划，尚非教材正文）：
% * 分式理想类群、非主理想及阶二理想类的算例
% * 理想正号约定下的 Weil 除子类群同构，与第10章 Picard 群及第20章类群定义相容
% * 对实数 s>1 定义 Dedekind ζ 函数，以有限素理想乘积和单调极限证明 Euler 乘积
% * 按有理素数归并因子，以通常收敛级数的 n 次幂给出统一上界
% <!-- 短节：不设小节 -->


理想分解修复了元素分解失效的问题，但仍留下一个具体问题：一个理想何时可以由单个元素生成？把主理想与一般理想的差别组织成群，便得到继续学习代数数论所需的一个基本对象。

\ProseParagraphBreak
沿用\cref{b3:def:fractional-invertible-ideal,b3:def:ideal-class-picard}。在数域 \(K\) 中，令 \(R=\mathcal O_K\)；非零分式理想可以等价地看作 \(K\) 的非零有限生成 \(R\)-子模，或可用一个非零 \(d\in R\) 清除分母的非零子模。乘法由有限和 \(\sum a_ib_i\) 定义；逆元为
\[
I^{-1}=\{x\in K:xI\subseteq R\}.
\]
\cref{b3:cor:dedekind-fractional-group}已经证明它们构成交换群，且每个分式理想唯一分解成素理想的整数次幂。理想类群是商群
\[
\operatorname{Cl}(R)=
\{\text{非零分式理想}\}/\{aR:a\in K^\times\}.
\]
同一理想类中的理想相差一个非零标量。群平凡恰好表示每个非零理想都主生成，见\cref{b3:prop:dedekind-picard-class}。

\ProseParagraphBreak
这里的 \(R\) 是一维 Dedekind 整环，非零素理想恰为高度一素理想。分式理想的唯一分解给出群同构
\[
I=\prod_{\mathfrak p}\mathfrak p^{a_{\mathfrak p}}
\longmapsto\sum_{\mathfrak p}a_{\mathfrak p}[\mathfrak p]
\in\operatorname{Div}(R),
\]
其中各 \(a_{\mathfrak p}\) 为整数且只有有限个非零；乘法变成除子的加法。主分式理想 \(aR\)、\(a\in K^\times\) 恰被送到主除子
\[
\operatorname{div}(a)=\sum_{\mathfrak p}v_{\mathfrak p}(a)[\mathfrak p].
\]
取商以后，本节的理想类群便与\cref{b3:def:weil-divisor-class-group}定义的 Weil 除子类群自然同构。这里与\cref{b3:thm:class-group-reflexive}采用同一理想正号约定：素理想 \(\mathfrak p\) 送到除子 \([\mathfrak p]\)。再由\cref{b3:prop:dedekind-picard-class}，把理想看作秩一投射模还得到 \(\operatorname{Cl}(R)\simeq\operatorname{Pic}(R)\)。

\begin{example}\label{b3:exm:nonprincipal-ideal-minus-six}
取 \(K=\mathbb Q(\sqrt{-6})\)、\(R=\mathbb Z[\sqrt{-6}]\)，令
\(\mathfrak p=(2,\sqrt{-6})\)。二次整基公式说明 \(R\) 确为整数环，商环为 \(\mathbb F_2\)，故 \(N(\mathfrak p)=2\)。若 \(\mathfrak p=(a+b\sqrt{-6})\)，主理想的指数等于乘法矩阵行列式的绝对值，因此
\[
2=|N_{K/\mathbb Q}(a+b\sqrt{-6})|=a^2+6b^2,
\]
这没有整数解。所以 \(\mathfrak p\) 非主。

\ProseParagraphBreak
另一方面，\(\mathfrak p^2\) 由
\(4,\,2\sqrt{-6},\,-6\) 生成，均属于 \(2R\)，而 \(-6+2\cdot4=2\)。故 \(\mathfrak p^2=2R\)。类群中 \([\mathfrak p]\) 的阶恰为 \(2\)。这个计算与第10章的 \(\mathbb Q(\sqrt{-5})\) 例子使用同一方法：先以剩余域确定理想范数，再用一个整数二次型排除主生成。它只证明存在这样的非平凡理想类，并没有枚举整个类群。
\end{example}

理想类群把“差多少才成为主理想”变成一个群问题。进一步的类域论把理想类及带同余条件的类似对象，与数域的交换 Galois 扩张\footnote{伽罗瓦（Évariste Galois，1811-1832），法国数学家，研究方程的根与置换群之间的联系，奠定了 Galois 理论的基础。}联系起来。带同余条件的理想类形成射线类群，无穷素点及分歧则使这些群与扩张的联系更加精细。理想分解、局部赋值和完备化为研究这些互反关系提供了算术基础。

\ProseParagraphBreak
另一种研究方法把全部理想范数放进一个级数。它从唯一分解出发，与素数分布和理想类的计数问题发生联系。下面证明 Euler 乘积\footnote{欧拉（Leonhard Euler，1707-1783），瑞士数学家，研究数论、分析与无穷级数，建立了把级数同素数乘积联系起来的经典公式。}对实数 \(s>1\) 的公式。

\begin{proposition}\label{b3:prop:dedekind-zeta-euler}
设 \(K\) 为数域，\(n=[K:\mathbb Q]\)，\(\mathcal O_K\) 为其代数整数环，并记非零理想 \(I\subseteq\mathcal O_K\) 的范数为 \(N(I)=|\mathcal O_K/I|\)。对实数 \(s>1\)，级数
\[
\zeta_K(s)=\sum_{0\ne I\subseteq\mathcal O_K}N(I)^{-s}
\]
收敛。其和称为 \(K\) 的\term[Dedekind-zeta-hanshu]{Dedekind \(\zeta\) 函数}[Dedekind zeta function]，并满足
\[
\zeta_K(s)=\prod_{\mathfrak p}
\bigl(1-N(\mathfrak p)^{-s}\bigr)^{-1},
\]
其中乘积遍历非零素理想。
\end{proposition}
\begin{proof}
理想分解、中国剩余定理及素理想幂的范数公式给出范数的乘法性。对于任意有限组非零素理想 \(S\)，展开有限个几何级数，得
\[
P_S=\prod_{\mathfrak p\in S}
\bigl(1-N(\mathfrak p)^{-s}\bigr)^{-1}
=\sum_{\substack{0\ne I\subseteq\mathcal O_K\\
\text{所有素因子均在 }S}}N(I)^{-s}.
\]
理想唯一分解保证每个理想恰计入一次，单位理想对应各指数均为零的项。

对固定的有理素数 \(p\)，\cref{b3:prop:prime-decomposition-degree}给出 \(p\) 上方至多有 \(n\) 个素理想，且每个剩余次数 \(f_{\mathfrak p}\ge1\)。于是
\[
\prod_{\mathfrak p\mid p}
\bigl(1-N(\mathfrak p)^{-s}\bigr)^{-1}
=\prod_{\mathfrak p\mid p}
\bigl(1-p^{-f_{\mathfrak p}s}\bigr)^{-1}
\le(1-p^{-s})^{-n}.
\]
令 \(T\) 为 \(S\) 中各素理想下方的有理素数所成的有限集，便有
\[
P_S\le\left(\prod_{p\in T}(1-p^{-s})^{-1}\right)^n
\le\left(\sum_{m\ge1}m^{-s}\right)^n<\infty.
\]
第二个不等式只需整数唯一分解：展开有限乘积，所得和遍历素因子均属于 \(T\) 的正整数，是右侧普通级数的一部分。普通级数在 \(s>1\) 时由积分比较收敛，因而此界不依赖 \(S\)。

令 \(S\) 依次包含前若干个有理素数上方的全部素理想。这些有限乘积单调有界，而每个非零理想只含有限多个素因子，最终都被计入。对非负项取单调极限，得到理想级数的收敛及所列 Euler 乘积等式。
\end{proof}

这个公式的代数内容已经全部体现在理想唯一分解中。若要研究 \(s=1\) 附近的行为，则需要额外的分析工具；若要由它计算类数，还需要单位群及数域的实、复嵌入。这些对象使理想的乘法结构与单位、嵌入及解析性质发生联系。


\begin{chaptersummary}
整数环的有限性来自迹配对：整元素的坐标受到一个非零整数行列式的控制。它的一维性与整闭性又使每个非零素理想处的局部环成为 DVR，各处的赋值由此组成非零理想的唯一分解。元素的唯一分解可以失效，理想的唯一分解仍然成立。

\ProseParagraphBreak
整基决定域判别式；若只计算某个子环的判别式，必须扣除指数的平方。模素数分解在指数互素的条件下给出素理想分解，分歧恰好由整数环判别式的素因子检测。整体 \(p\)-进完备化 \(\mathcal O_K\otimes_{\mathbb Z}\mathbb Z_p\) 分裂成 \(p\) 上方全部素理想对应的局部完备因子；选定一个 \(\mathfrak p\mid p\)，才得到其中一个完备整数环及其分式域。

\ProseParagraphBreak
理想类群把理想的乘法改写成类的加法；在一维整数环中，它也就是 Weil 除子类群，并与 Picard 群自然同构。Dedekind ζ 函数按范数对所有非零理想求和，理想唯一分解则把这次求和改写成素理想乘积。在实数 \(s>1\) 时，按有理素数归并各因子，便能由通常的 ζ 级数控制其收敛。
\end{chaptersummary}

\BookExercises

\begin{exercise}\label{b3:ex:A-clear-denominator}
令 \(K=\mathbb Q(\sqrt5)\)、\(a=(1+\sqrt5)/6\)。求使 \(da\) 为代数整数的最小正整数 \(d\)，写出其首一整系数方程，并说明坐标的共同分母 \(6\) 为什么不一定给出最小的 \(d\)。
\end{exercise}
\begin{solution}
若 \(da\) 整，其迹 \(d/3\) 必须为整数，故 \(3\mid d\)。取 \(d=3\) 时，\(3a=(1+\sqrt5)/2\) 满足 \(T^2-T-1=0\)，确为代数整数。因此最小值是 \(3\)。有理基 \(1,\sqrt5\) 的坐标分母只给出一个足够大的清分母数；真正的整数环还包含半整数坐标的整元素。
\end{solution}

\begin{exercise}\label{b3:ex:A-trace-bound}
在 \(K=\mathbb Q(\sqrt2)\) 中，对整元素基 \(1,\sqrt2\) 计算迹矩阵与行列式，并写出迹配对证明给出的格包含。解释这个包含为何通常不是最精确的整数环描述。
\end{exercise}
\begin{solution}
迹矩阵为 \(\operatorname{diag}(2,4)\)，行列式为 \(8\)。因此令 \(L=\mathbb Z+\mathbb Z\sqrt2\)，一般论证给出
\[
L\subseteq\mathcal O_K\subseteq\dfrac18L.
\]
二次整基公式实际上给出 \(\mathcal O_K=L\)。伴随矩阵只给一个统一的分母界，并未利用每个候选元素的全部整性条件，故不要求等号成立于右侧。
\end{solution}

\begin{exercise}\label{b3:ex:A-quadratic-bases}
分别求 \(\mathbb Q(\sqrt{-3})\)、\(\mathbb Q(\sqrt2)\)、\(\mathbb Q(\sqrt{13})\) 的整基与域判别式，并写出所选非有理整基元素的最小多项式。
\end{exercise}
\begin{solution}
整基依次可取 \(1,(1+\sqrt{-3})/2\)，\(1,\sqrt2\)，\(1,(1+\sqrt{13})/2\)。相应判别式为 \(-3,8,13\)，最小多项式为 \(T^2-T+1\)、\(T^2-2\)、\(T^2-T-3\)。这些结论均由无平方因子整数模 \(4\) 的分类得到。
\end{solution}

\begin{exercise}\label{b3:ex:A-index-discriminant}
证明 \([\mathbb Z[(1+\sqrt5)/2]:\mathbb Z[\sqrt5]]=2\)，并直接检验两个判别式之间的指数平方公式。
\end{exercise}
\begin{solution}
写 \(\omega=(1+\sqrt5)/2\)，则 \(\sqrt5=2\omega-1\)。基 \(1,\sqrt5\) 相对于 \(1,\omega\) 的坐标矩阵为
\(\left(\begin{smallmatrix}1&-1\\0&2\end{smallmatrix}\right)\)，行列式为 \(2\)，所以指数为 \(2\)。两个判别式分别为 \(20,5\)，确有 \(20=2^2\cdot5\)。
\end{solution}

\begin{exercise}\label{b3:ex:A-gaussian-primes}
在 \(\mathbb Z[i]\) 中分解理想 \((2)\)、\((3)\)、\((5)\)，给出各素因子的分歧指数与剩余次数。
\end{exercise}
\begin{solution}
整数环等于 \(\mathbb Z[i]\)，使用 \(T^2+1\)。模 \(2\) 为 \((T+1)^2\)，故
\((2)=(1+i)^2\)，因为 \((1+i)^2=2i\) 且 \(i\) 是单位；分歧指数为 \(2\)、剩余次数为 \(1\)。模 \(3\) 不可约，所以 \((3)\) 为素理想，指数 \(1\)、剩余次数 \(2\)。模 \(5\) 的根为 \(2,-2\)，故
\((5)=(5,i-2)(5,i+2)\)，两因子的指数与剩余次数均为 \(1\)；也可写成主理想 \((2+i)(2-i)\)。
\end{solution}

\begin{exercise}\label{b3:ex:A-ideal-containment}
设 \(K\) 为数域，\(R=\mathcal O_K\) 为其整数环，非零理想 \(I,J\) 的素理想指数分别为 \(a_{\mathfrak p},b_{\mathfrak p}\)，缺失的指数记为零。证明 \(I\subseteq J\) 当且仅当所有 \(a_{\mathfrak p}\ge b_{\mathfrak p}\)，并求 \(I+J\) 与 \(I\cap J\) 的指数。
\end{exercise}
\begin{solution}
在每个 离散赋值环 \(R_{\mathfrak p}\) 中，统一化元的幂所生成的理想的包含关系与指数大小相反。局部检测给出全局的包含判据。局部和取较小指数，局部交取较大指数；有限交与局部化相容，再用局部检测，得到
\[
v_{\mathfrak p}(I+J)=\min(a_{\mathfrak p},b_{\mathfrak p}),\qquad
v_{\mathfrak p}(I\cap J)=\max(a_{\mathfrak p},b_{\mathfrak p}).
\]
\end{solution}

\begin{exercise}\label{b3:ex:A-principal-norm}
令 \(R=\mathbb Z[\sqrt{-6}]\)、\(a=2+\sqrt{-6}\)。构造同构 \(R/(a)\cong\mathbb Z/10\mathbb Z\)，并证明
\[
(a)=(2,\sqrt{-6})(5,\sqrt{-6}+2).
\]
用主理想范数公式独立核对商环的基数。
\end{exercise}
\begin{solution}
记 \(u=\sqrt{-6}\)。商去 \(u+2\) 后有 \(u=-2\)，关系 \(u^2+6=0\) 变成 \(10=0\)，因而代入 \(u\mapsto-2\) 给出所述同构。两个理想的商分别为 \(\mathbb F_2,\mathbb F_5\)，故均为素理想，互素且范数为 \(2,5\)。元素 \(a\) 属于两者，所以 \((a)\) 包含于它们的交，也就是乘积。另一方面，\cref{b3:lem:principal-ideal-norm}给出 \(N((a))=2^2+6=10\)，乘积的范数也是 \(10\)。包含的两个有限指数理想具有相同指数，故相等。
\end{solution}

\begin{exercise}\label{b3:ex:A-discriminant-primes}
令 \(K=\mathbb Q(\sqrt{-3})\)、\(\eta=(1+\sqrt{-3})/2\)。证明 \(5,7\) 在 \(K\) 中都不分歧，但前者惰性、后者完全分裂；写出素理想分解及各因子的 \((e,f)\)。这说明判别式能检测分歧，却不能单独决定所有未分歧素数的分解型。
\end{exercise}
\begin{solution}
由整基公式，\(R=\mathbb Z[\eta]\)，域判别式为 \(-3\)，所以 \(5,7\) 都不分歧。最小多项式为 \(f(T)=T^2-T+1\)。模 \(5\) 的判别式是 \(2\)，不是平方，故 \(f\) 不可约；于是 \(5R\) 是素理想，\((e,f)=(1,2)\)。模 \(7\) 有 \(f=(T-3)(T-5)\)，所以
\[
7R=(7,\eta-3)(7,\eta-5),
\]
两个因子均有 \((e,f)=(1,1)\)。两处都没有重复因子，但剩余域次数不同。
\end{solution}

\begin{optionalexercise}\label{b3:ex:A-eleven-adic-lifts}
设 \(a_1,a_2\in\mathbb Z_{11}\) 为 \(T^2-T-1\) 的两个根，分别同余于 \(4,8\pmod{11}\)。从已知的 \(37,85\pmod{121}\) 继续计算它们模 \(1331\) 的值，并写出求值映射
\[
\mathbb Z_{11}[W]/(W^2-W-1)\longrightarrow\mathbb Z_{11}^2,
\qquad h\longmapsto(h(a_1),h(a_2))
\]
的显式逆映射。
\end{optionalexercise}
\begin{solution}
记 \(f(T)=T^2-T-1\)。有 \(f(37)=1331\)，故第一根无需改变下一位，仍为 \(37\pmod{1331}\)。另一根写为 \(85+121c\)，因 \(f(85)=121\cdot59\)、\(f'(85)=169\)，须解 \(59+169c\equiv0\pmod{11}\)，得 \(c=10\)，所以第二根为 \(1295\pmod{1331}\)。二者之和仍同余于 \(1\)。由于 \(a_2-a_1\equiv4\pmod{11}\)，这个差在 \(\mathbb Z_{11}\) 中为单位。给定 \((r,s)\)，令
\[
v=\dfrac{s-r}{a_2-a_1},\qquad
u=\dfrac{a_2r-a_1s}{a_2-a_1}.
\]
则 \(u+vW\) 在两根处取值为 \(r,s\)。每个商类唯一写成 \(u+vW\)，故这个公式给出所求逆映射。
\end{solution}

\begin{exercise}\label{b3:ex:A-local-product}
写出 \(\mathbb Q(i)\otimes_{\mathbb Q}\mathbb Q_p\) 在 \(p=2,3,5\) 时的因子数及各因子的次数、分歧指数与剩余次数。
\end{exercise}
\begin{solution}
在 \(p=2\) 时有一个次数二因子，\((e,f)=(2,1)\)；在 \(p=3\) 时有一个次数二因子，\((e,f)=(1,2)\)；在 \(p=5\) 时有两个次数一因子，均为 \(\mathbb Q_5\)，各有 \((e,f)=(1,1)\)。这是素理想分解与完备化乘积公式的直接应用，且每种情形的次数之和均为 \(2\)。
\end{solution}

\begin{exercise}\label{b3:ex:A-fractional-valuations}
设 \(K\) 为数域，\(R=\mathcal O_K\) 为其整数环。证明每个非零分式理想唯一写成有限乘积
\(\prod_{\mathfrak p}\mathfrak p^{a_{\mathfrak p}}\)，其中 \(a_{\mathfrak p}\in\mathbb Z\)。据此证明 \(R\) 的理想类群平凡当且仅当 \(R\) 为 PID。
\end{exercise}
\begin{solution}
选 \(0\ne d\in R\) 使 \(dI\subseteq R\)。分别分解积分理想 \(dI,dR\)，再用已证的分式理想群取商，即得整数指数分解。唯一性由每个 DVR 中分式理想中统一化元的幂次得到。如果类群平凡，任何非零积分理想均为某个 \(aR\)，且 \(a\in aR\subseteq R\)，所以为主理想；零理想也主，故是 PID。反过来，PID 中先把分式理想清分母再写成主理想，便知每个分式理想均主。
\end{solution}

\begin{exercise}\label{b3:ex:A-two-factorizations}
在 \(\mathbb Z[\sqrt{-5}]\) 中证明 \(2,3,1+\sqrt{-5},1-\sqrt{-5}\) 都不可约，且
\[
6=2\cdot3=(1+\sqrt{-5})(1-\sqrt{-5})
\]
是两种不等价的不可约分解。
\end{exercise}
\begin{solution}
元素范数为 \(a^2+5b^2\)，非零非单位范数大于 \(1\)。范数 \(2,3\) 均不能表示；故范数分别为 \(4,9\) 的 \(2,3\) 不可能分成两个非单位。范数为 \(6\) 的 \(1\pm\sqrt{-5}\) 也不能分解，因为非平凡整数因子只能为 \(2,3\)。单位只有 \(\pm1\)；四个元素按上述配对均不相伴，故两种分解不等价。这不违背非零理想的唯一分解。
\end{solution}

\begin{exercise}\label{b3:ex:A-euler-factors}
设 \(K\) 为数域，\(R=\mathcal O_K\)，\(p\) 为有理素数。记 \(p\) 上方剩余次数为 \(1,2\) 的素理想数分别为 \(r_1,r_2\)。证明范数为 \(p\) 与 \(p^2\) 的积分理想数依次为
\[
r_1,\qquad r_2+\dfrac{r_1(r_1+1)}2.
\]
再计算 \(K=\mathbb Q(\sqrt5)\)、\(p=2,5,11\) 时这两个理想数，说明为何不能把分歧指数直接当成计数的重数。
\end{exercise}
\begin{solution}
范数为 \(p\) 的理想只能是剩余次数一的某个素理想。范数为 \(p^2\) 的理想则是剩余次数二的一个素理想，或两个剩余次数一的素理想的乘积，允许二者相同。理想唯一分解使后者按无序可重复二元组计数，故共有 \(r_1(r_1+1)/2\) 个。这也分别是形式乘积 \(\prod_{\mathfrak p\mid p}(1-T^{f_{\mathfrak p}})^{-1}\) 的 \(T,T^2\) 系数。对 \(\mathbb Q(\sqrt5)\)，三个素数处的两项计数依次为 \((0,1)\)、\((1,1)\)、\((2,3)\)。在 \(5\) 处只有一个范数 \(5\) 的素理想，\(e=2\) 并没有产生第二个这样的理想。
\end{solution}
